\documentclass[11pt]{article}
\usepackage[fleqn]{amsmath}
\setlength{\mathindent}{1em}
\usepackage{amssymb, amsthm}
\usepackage[utf8]{inputenc}
\usepackage[margin=1in]{geometry}
\usepackage{graphicx}
\usepackage{enumitem}
\theoremstyle{plain}
\newtheorem{proposition}{Proposition}
\newtheorem*{proposition*}{Proposition}
\newcommand{\fit}[1]{\resizebox{\ifdim\width>\linewidth\linewidth\else\width\fi}{!}{#1}}
\setlength{\parindent}{0pt}
\begin{document}
\begin{proposition*}[cauchy-convex-closed]
Suppose $U \text{ s.t. } \operatorname{is\text{-}convex\text{-}cc}(U)$. Let $f \text{ be s.t. } \operatorname{holomorphic\text{-}on}(U, f)$. Let $\gamma$, $\gamma'$ and $a$ be arbitrary. Let $b \text{ be s.t. } \operatorname{is\text{-}road}(\gamma, \gamma', a, b)$. If $\operatorname{trace}(\gamma, a, b) \subseteq U$, then if $\gamma(b) = \gamma(a)$ then  $\operatorname{line\text{-}int}(f, \gamma, \gamma', a, b) = 0$
\end{proposition*}

\begin{proof}
{\small\emph{The numbers in the margin are step numbers of the full proof (\texttt{proof-tex}); a range such as 13--15 is several steps shown as one row.  A claim is proved by the indented rows under it.}}\par\smallskip
\begin{itemize}[leftmargin=2.8em, itemsep=4pt, parsep=1pt]
  \item[\textbf{8}] Clearly $\operatorname{is\text{-}open}(\operatorname{nf\text{-}metric\text{-}space}(\operatorname{cc\text{-}normed\text{-}field}), U)$.
  \item[\textbf{9}] By the proposition \emph{cc-open-subset-cc} at $u := U$ (its hypothesis $\operatorname{is\text{-}open}(\operatorname{nf\text{-}metric\text{-}space}(\operatorname{cc\text{-}normed\text{-}field}), U)$ is in context): $U \subseteq \mathbb{C}$
  \item[\textbf{10--13}] Clearly
\begin{equation*}
\begin{array}{@{}l@{}}
f \in (U \to \mathbb{C}) \\
\begin{array}{@{}l@{}} \forall y \in U. \\ \quad \operatorname{is\text{-}continuous\text{-}at}(\operatorname{subspace\text{-}ms}(\operatorname{nf\text{-}metric\text{-}space}(\operatorname{cc\text{-}normed\text{-}field}), U), \\ \quad \quad \operatorname{nf\text{-}metric\text{-}space}(\operatorname{cc\text{-}normed\text{-}field}), f, y) \end{array} \\
\operatorname{is\text{-}path}(\gamma, a, b) \\
\gamma \in (\operatorname{ccint}(a, b) \to \mathbb{C})
\end{array}
\end{equation*}

  \item[\textbf{14}] By the proposition \emph{is-path-endpoints} at $\gamma := \gamma$, $a := a$, $b := b$ (its hypothesis $\operatorname{is\text{-}path}(\gamma, a, b)$ is in context): $(a \in \mathbb{R}) \wedge ((b \in \mathbb{R}) \wedge (a < b))$
  \item[\textbf{15--17}] Hence $a \in \mathbb{R}$, $a < b$ and $b \in \mathbb{R}$. \textbf{Claim 1} (proved in steps 18--23): $a \in \operatorname{ccint}(a, b)$
  \begin{itemize}[leftmargin=2.8em, itemsep=4pt, parsep=1pt]
  \item[\textbf{18--23}] By the proposition \emph{ccint-membership}, reduce to $(a \in \mathbb{R}) \wedge ((a \leq a) \wedge (a \leq b))$  \emph{(it holds by assumption; this proves Claim 1)}
  \end{itemize}
  \item[\textbf{24--25}] Clearly $\gamma(a) \in \operatorname{trace}(\gamma, a, b)$ and $\gamma(a) \in U$.
  \item[\textbf{26}] By the proposition \emph{cauchy-convex-primitive} at $u := U$, $f := f$, $o := \gamma(a)$ (its hypotheses $\gamma(a) \in U$, $\operatorname{holomorphic\text{-}on}(U, f)$ and $\operatorname{is\text{-}convex\text{-}cc}(U)$ are in context):
\begin{equation*}
\begin{array}{@{}l@{}} \exists g \in (U \to \mathbb{C}). \\ \quad \operatorname{holomorphic\text{-}on}(U, g) \wedge \\ \quad \forall z \in U.\; \operatorname{is\text{-}diff\text{-}on}(\operatorname{cc\text{-}normed\text{-}field}, U, g, z, f(z)) \end{array}
\end{equation*}

  \item[\textbf{27--31}] Hence $g \in (U \to \mathbb{C})$, $\forall z \in U.\; \operatorname{is\text{-}diff\text{-}on}(\operatorname{cc\text{-}normed\text{-}field}, U, g, z, f(z))$ and $\operatorname{holomorphic\text{-}on}(U, g)$. By the proposition \emph{line-int-closed-road-of-derivative} at $u := U$, $f := g$, $d := f$, $\gamma := \gamma$, $\gamma' := \gamma'$, $a := a$, $b := b$ (its hypotheses $\gamma(b) = \gamma(a)$, $\operatorname{trace}(\gamma, a, b) \subseteq U$, $\operatorname{is\text{-}road}(\gamma, \gamma', a, b)$ and $f \in (U \to \mathbb{C})$ are in context): $\operatorname{line\text{-}int}(f, \gamma, \gamma', a, b) = 0$  \emph{(holds by assumption)}
\end{itemize}
\end{proof}
\end{document}
